\documentclass[12pt,a4paper]{article}
\usepackage[russian]{babel}
\usepackage{fontspec}
\usepackage{geometry}
\usepackage{booktabs}
\usepackage{longtable}
\usepackage{array}
\usepackage{multirow}
\usepackage{enumitem}
\usepackage{titlesec}
\usepackage{xcolor}
\usepackage{siunitx}

\usepackage{pdfpages}

\usepackage{tablefootnote}

% Дополнительные пакеты для данного документа
\usepackage{caption}
\usepackage{subcaption}
\usepackage{listings}
% \usepackage{hyperref}

\geometry{a4paper, margin=25mm}
\setmainfont{Liberation Serif}
\setsansfont{Liberation Sans}

\definecolor{adblue}{RGB}{0,84,166}

\titleformat{\section}{\large\bfseries\color{adblue}}{}{0em}{}[\titlerule]
\titleformat{\subsection}{\normalsize\bfseries}{}{0em}{}
\titleformat{\subsubsection}{\normalsize\itshape}{}{0em}{}

\sisetup{output-decimal-marker={,}}

\usepackage{tikz}
\usetikzlibrary{positioning, calc, backgrounds, math, 3d, perspective, arrows.meta, graphs, graphdrawing, fit, shapes.geometric, trees}
    

\usepackage{pgfplots} % для построения графиков
\pgfplotsset{compat=1.18} % версия для совместимости
\usepackage{tkz-euclide}

% --- Подключение пакета алгоритмов (без флага algo2e, чтобы вернуть \begin{algorithm}) ---
\usepackage[ruled,vlined]{algorithm2e}
%\usepackage{caption}
\newcommand{\Gray}[1]{\textcolor{gray}{#1}} 
% Настройка подписей
\DeclareCaptionLabelFormat{gost}{#1 #2}
\captionsetup[table]{labelformat=gost,labelsep=endash,justification=justified,singlelinecheck=false,font=normal}
\captionsetup[figure]{labelformat=gost,labelsep=endash,justification=centering,font=normal}

% Локализация algorithm2e на русский (ОБЯЗАТЕЛЬНО после загрузки пакета)
\SetAlgorithmName{Алгоритм}{}{}
\SetAlgoCaptionSeparator{---}
\SetKwInput{Input}{Вход}
\SetKwInput{Output}{Выход}

\usepackage{url}
% Говорим пакету использовать моноширинный шрифт (как у texttt)
\Urlmuskip=0mu plus 1mu % гибкие пробелы для формул, если нужно
\DeclareUrlCommand\code{\urlstyle{tt}} 

% Библиотека схемотехники
\usepackage[siunitx, RPvoltages, american]{circuitikzgit}
\ctikzset{
    amplifiers/fill=cyan,
    sources/fill=green!20,
    diodes/fill=white,
    resistors/fill=violet,
    grounds/scale=1.2,
}

\usepackage{iftex}

\ifXeTeX % Если используется TeTeX или TeLaTeX
  \typeout{Режим XeTeX}
  \usepackage{xltxtra}
  \usepackage{fontspec}
  \usepackage{polyglossia}
  \setmainlanguage{russian}
  \setotherlanguage{english}
  \setkeys{russian}{babelshorthands=true} % По идее включает <<, >>, -- ---

  \setmainfont{Times New Roman}[Ligatures=TeX]
  \setromanfont{Times New Roman}[Ligatures=TeX]
  \setsansfont{Arial}[Ligatures=TeX]
  \setmonofont{Courier New}[Ligatures=TeX] 

  \newfontfamily{\cyrillicfont}{Times New Roman}
  \newfontfamily{\cyrillicfontrm}{Times New Roman}
  \newfontfamily{\cyrillicfonttt}{Courier New}
  \newfontfamily{\cyrillicfontsf}{Arial}
\fi

\ifLuaTeX % Если используется LuaLatex
  \typeout{Режим LuaTeX}
%  \usepackage{luatextra}
  \usepackage{fontspec}
  \defaultfontfeatures{Ligatures={TeX}}
  \usepackage{polyglossia}
  \setmainlanguage{russian}
  \setotherlanguage{english}
  \setmainfont{Times New Roman}[Script=Cyrillic, Ligatures=TeX]
  \setromanfont{Times New Roman}[Script=Cyrillic, Ligatures=TeX]
  \setsansfont{Arial}[Script=Cyrillic, Ligatures=TeX]
  \setmonofont{Courier New}[Ligatures=TeX]
  \newfontfamily{\cyrillicfont}{Times New Roman}[Script=Cyrillic, Ligatures=TeX]
  \newfontfamily{\cyrillicfontrm}{Times New Roman}[Script=Cyrillic, Ligatures=TeX]
  \newfontfamily{\cyrillicfonttt}{Courier New}
  \newfontfamily{\cyrillicfontsf}{Arial}[Script=Cyrillic, Ligatures=TeX]

  \usegdlibrary{trees}
  \usegdlibrary{force}
\fi

\ifPDFTeX % Если используется 8битный PDFTex
  \typeout{8-битный режим pdfTex}
  % For pdfTeX we must use old
  % 8-bit font technologies
  \usepackage[T2A]{fontenc}
  \usepackage[english,russian]{babel}
  \usepackage[utf8]{inputenc}
  \usepackage[english, russian]{babel}

  %\usepackage{newtxmath}
  \usepackage{amsmath}
  \usepackage{amsthm}
  \usepackage{amssymb}
  \usepackage{amsfonts}
  \usepackage{mathtext}
\fi

\usepackage{amsthm}  % Возможности AMSMath
\usepackage{amsmath}
\usepackage{amssymb}

\usepackage{esvect} % Для красивых стрелок
\usepackage{bm} % Для жирных символов

%Hyphenation rules
\usepackage{hyphenat}
\hyphenation{ма-те-ма-ти-ка вос-ста-нав-ли-вать ото-бра-жа-ют-ся сге-не-ри-ро-ван-ном сим-во-лы}
% \usepackage[english, russian]{babel}

\usepackage{graphicx}
\graphicspath{ {./images/} }

\usepackage{empheq}

\usepackage{float}

% \usepackage{comment}

\usepackage[hidelinks,
    unicode=true,          % поддерживает Unicode в закладках
    psdextra=true,      % автоматически конвертирует простую математику в текст
    % focusdefault=false
]{hyperref}


% Окружение "Теорема"
\newtheorem{theorem}{Теорема}[section]
\newtheorem{corollary}{Corollary}[theorem]
\newtheorem{lemma}[theorem]{Lemma}

% Окружение "Определение"
\theoremstyle{definition} % Прямой шрифт, нумерация
\newtheorem{definition}{Определение}[section]
\newtheorem{example}{Пример}[section]

% Окружение "Замечание"
\theoremstyle{remark} % Прямой шрифт, без нумерации
% Окуржение "Решение"
\newtheorem*{solution}{Решение}
\newtheorem*{remark}{Замечание}

\usepackage{polynom}
\usepackage{tabularx}
\usepackage{cancel}

\addto\captionsrussian{%
  \renewcommand{\figurename}{Рис.}%
  \renewcommand{\tablename}{Табл.}%
}

\renewcommand{\arraystretch}{1.5} % Высота строки таблицы

\renewcommand{\labelitemi}{\(\circ\)} % Бюллет -- пустой кружок в itemize


\begin{document}

\begin{algorithm}[H]
% \SetCustomAlgoRuKeywords
\DontPrintSemicolon
\SetKwInOut{KwIn}{Вход}
\SetKwInOut{KwOut}{Выход}

\TitleOfAlgo{Сквозной конвейер восстановления и апскейла RAW-изображений на базе NAFNet}

\KwIn{Каталог исходных изображений высокого разрешения $D_{\text{source}}$, конфигурационный файл параметров \texttt{config.yml}, количество эпох обучения $T$, размер батча $B$, целевые масштабы патчей $S_{\text{lq}}, S_{\text{gt}}$, коэффициент апскейла $F_{\text{scale}}$.}
\KwOut{Оптимизированные веса нейросетевой модели $\Theta_{\text{opt}}$, файлы сквозной аналитики \texttt{train\_metrics.csv}, \texttt{val\_metrics.csv}, верификационные превью $\text{Img}_{\text{lq}}, \text{Img}_{\text{gt}}$.}

\BlankLine
\Gray{/* Этап I: Инициализация и верификация дисковой подсистемы */}
Парсинг и валидация аргументов командной строки через модуль \texttt{config.py}\;
Инициализация двухпоточного логгера \texttt{setup\_logger()} в файл \texttt{pipeline.log}\;
\If{выставлен флаг \texttt{--clean-training}}{
    Выборочное удаление старых каталогов \texttt{checkpoints}, \texttt{tb\_logs}, \texttt{val\_predictions}\;
}

\BlankLine
\Gray{/* Этап II: Синхронная генерация синтетического датасета */}
\If{каталог датасета пуст {\bf или} выставлен флаг \texttt{--clean-dataset}}{
    Очистка директорий \texttt{train} и \texttt{test} внутри корня $D_{\text{root}}$\;
    Разбиение файлов из $D_{\text{source}}$ на подмножества $\text{Files}_{\text{train}}$ и $\text{Files}_{\text{test}}$ в пропорции $\text{ratio}$\;
    \ForEach{изображение $I_{\text{rgb}}$ в $\text{Files}_{\text{train}}, \text{Files}_{\text{test}}$}{
        Чтение матрицы пикселей и приведение к 8-битному RGB-пространству\;
        Вырезание центрального патча $I_{\text{crop}}$ размера $S_{\text{gt}} \times S_{\text{gt}}$ с соблюдением чётности Bayer-фазы\;
        Применение PSF-размытия: $I_{\text{blur}} = \mathcal{Filter2D}(I_{\text{crop}}, \text{Kernel}_{\text{PSF}})$\;
        Накладывание маски мозаики фильтров: $I_{\text{bayer}} = \mathcal{ApplyBayerMask}(I_{\text{blur}}, \text{'RGGB'})$\;
        Генерация аддитивно-мультипликативного шума матрицы на основе SNR в дБ\;
        Упаковка 2D-мозаики в 4-канальный массив субканалов: $I_{\text{packed}} = \mathcal{ExtractSubchannels}(I_{\text{bayer}})$\;
        Сохранение пар на диск: $I_{\text{packed}} \rightarrow \text{uint16 tiff}$, $I_{\text{crop}} \rightarrow \text{uint8 png}$\;
    }
}

\BlankLine
\Gray{/* Этап III: Диагностический визуальный контроль данных */}
Извлечение случайного верификационного сэмпла из загруженного класса \texttt{CustomNEFPairDataset}\;
Применение быстрой билинейной интерполяции к 4-канальному входу\;
Конвертация цветового пространства из RGB в BGR: $I_{\text{bgr}} = \mathcal{Cv2Color}(I_{\text{rgb}}, \text{RGB2BGR})$\;
Запись диагностических превью $\text{Img}_{\text{lq}}$ и $\text{Img}_{\text{gt}}$ на диск через \texttt{cv2.imwrite()}\;

\BlankLine
\Gray{/* Этап IV: Нагрузочный Smoke-тест вычислительного графа */}
Инициализация базовой модели \texttt{NAFNet} и обёртки \texttt{NAFNetDemosaicSuperResolutionWrapper}\;
Перенос графа модели и тестового батча на ускоритель \texttt{torch.device('cuda')}\;
Прямой проход: $\hat{Y}_{\text{smoke}} = \text{Model}(X_{\text{smoke}})$\;
Обсчёт комбинированного критерия потерь: $L = L_1(\hat{Y}_{\text{smoke}}, Y_{\text{smoke}}) + L_{\text{FFL\_Log}}(\hat{Y}_{\text{smoke}}, Y_{\text{smoke}})$\;
Обратный проход градиентов $L\text{.backward()}$ и шаг оптимизатора AdamW\;

\BlankLine
\Gray{/* Этап V: Основной цикл оптимизации параметров (Обучение) */}
Инициализация логгеров TensorBoard и CSV-аналитики\;
\For{эпоха $e = 0$ \KwTo $T-1$}{
    Активация режима обучения: $\text{Model.train()}$\;
    \ForEach{батч $(X_{\text{lq}}, Y_{\text{gt}})$ в \texttt{DataLoader}}{
        Обнуление градиентов оптимизатора: $\text{optimizer.zero\_grad()}$\;
        Прямой проход через обёртку в полном разрешении: $\hat{Y} = \text{Model}(X_{\text{lq}})$\;
        Вычисление пространственно-частотного лосса: $L_{\text{total}}, l_1, l_{\text{ffl}} = \text{CombinedLoss}(\hat{Y}, Y_{\text{gt}})$\;
        Обратный проход градиентов: $L_{\text{total}}\text{.backward()}$\;
        Ограничение нормы градиентов: $\mathcal{ClipGradNorm}(\text{Model.parameters()}, 1.0)$\;
        Выполнение шага оптимизации: $\text{optimizer.step()}$\;
        \If{$e \pmod{\text{file\_log\_freq}} == 0$}{
            Расчёт дисперсии градиентов весов $\text{Grad}_{\text{var}}$ и коэффициента вариации $\text{TV}_{\text{ratio}}$\;
            Запись метрик в \texttt{train\_metrics.csv} и \texttt{train\_progress.log}\;
        }
    }
    Обновление скорости обучения шедулером: $\text{scheduler.step()}$\;
    \If{эпоха $e$ удовлетворяет \texttt{validation\_freq}}{
        Активация режима валидации: $\text{Model.eval()}$\;
        Расчёт метрик структурного сходства SSIM на тестовом Center Crop патче\;
        Запись результатов в \texttt{val\_metrics.csv}\;
    }
    \If{эпоха $e$ удовлетворяет \texttt{save\_checkpoint\_epoch}}{
        Сохранение состояний модели, оптимизатора и шедулера в \texttt{.pth} файл чекпоинта\;
    }
}
\Return{Оптимизированные веса модели $\Theta_{\text{opt}}$}
\caption{Алгоритм работы распределённого ноосферного конвейера RAW-restoration}
\end{algorithm}

\end{document}